\documentclass[10pt,a4paper]{article}
\usepackage[margin=18mm,headheight=14pt]{geometry}
\usepackage{fontspec}
\setmainfont{lmroman10-regular.otf}[BoldFont=lmroman10-bold.otf,ItalicFont=lmroman10-italic.otf,BoldItalicFont=lmroman10-bolditalic.otf]
\setsansfont{lmsans10-regular.otf}[BoldFont=lmsans10-bold.otf]
\setmonofont{lmmonolt10-regular.otf}[BoldFont=lmmonolt10-bold.otf,ItalicFont=lmmonolt10-oblique.otf]
\usepackage[ngerman]{babel}
\usepackage{amsmath,amssymb,booktabs,tabularx,array,fancyhdr,enumitem}
\usepackage[hidelinks]{hyperref}
\usepackage{tikz}
\usetikzlibrary{shapes.gates.logic.US,positioning,calc}
\tikzset{gate/.style={draw,minimum width=15mm,minimum height=10mm,font=\sffamily\small},wire/.style={thick}}
\pagestyle{fancy}\fancyhf{}
\fancyhead[L]{\sffamily AI-Hardware · Tag 5}
\fancyhead[R]{\sffamily Abschluss · Aufgabenblatt}
\fancyfoot[L]{\small Logikgatter und Boolesche Algebra}
\fancyfoot[R]{\thepage\ / 5}
\setlength{\emergencystretch}{3em}\setlength{\parindent}{0pt}\setlength{\parskip}{5pt}
\setlength{\tabcolsep}{7pt}\renewcommand{\arraystretch}{1.18}
\newcommand{\titlepagepart}[2]{\vspace*{-4mm}{\sffamily\small #1}\par{\sffamily\LARGE\bfseries #2}\par\vspace{2mm}}
\newcommand{\block}[1]{\par\vspace{2mm}{\sffamily\large\bfseries #1}\par}
\newcommand{\code}[1]{\texttt{#1}}
\newcommand{\antwort}[1]{\par\vspace{#1}\hrule\vspace{3mm}}
\newcommand{\N}{\mathrm{N}}
\newcommand{\lin}{\rule{11mm}{0.3pt}}
\begin{document}

\titlepagepart{01 / BOOLESCHE ALGEBRA · CA. 12 MINUTEN}{Tag 5 · Abschluss}
Name: \rule{65mm}{0.3pt}\hfill Datum: \rule{30mm}{0.3pt}

Dieses Blatt schließt Tag 5 ab. Es setzt die Zusammenfassung
\code{Tag05\_Zusammenfassung.pdf} voraus. Plane etwa 60 Minuten: rund
40 Minuten Papier (Seiten 1--4), rund 20 Minuten Logisim und Dokumentation
(Seite 5). Schreibe die Zwischenschritte auf, nicht nur das Ergebnis, und
notiere hinter jedem Umformungsschritt das verwendete Gesetz. Schlage erst
nach, wenn du eine Aufgabe ohne Hilfe versucht hast, und markiere sie dann.

\block{Stand im Ordner \code{01\_digital/02\_logikgatter/tag05}}
Vorhanden sind die achtseitige Zusammenfassung und vier Logisim-Dateien
(\code{NandToXor}, \code{NandToNOR}, \code{DeMorgan1\_NOR},
\code{DeMorgan2\_NAND}). Damit sind XOR und NOR aus NAND sowie beide
De-Morgan-Vergleiche aufgebaut. Offen für die Erfolgskontrolle sind NOT, AND
und OR als eigene geprüfte Aufbauten, die im Plan vorgesehene Sammeldatei
\code{tag05\_nand.circ}, die schriftlich festgehaltenen Wahrheitstabellen und
der Commit. Darauf zielt Seite 5.

\block{1. Vereinfachen mit den Gesetzen}
\textbf{a)} Vereinfache $Y=\overline{A}B+AB+A\overline{B}$ schrittweise bis zur
kürzesten Form. Gib pro Zeile das benutzte Gesetz an.
\antwort{26mm}
\textbf{b)} Zeige $Y=(A+B)(A+\overline{B})=A$. Verwende das Distributivgesetz
und anschließend das Komplement- und das Neutralitätsgesetz.
\antwort{24mm}
\textbf{c)} Vereinfache $Y=\overline{\overline{A}+\overline{B}}$. Benenne beide
Regeln, die du dafür brauchst.
\antwort{20mm}
\textbf{d)} Streiche im Konsensterm $Y=AB+\overline{A}C+BC$ den redundanten
Term. Begründe in einem Satz, warum er nichts hinzufügt.
\antwort{24mm}

\newpage
\titlepagepart{02 / DE MORGAN · CA. 12 MINUTEN}{Die Negation nach innen ziehen}

\block{2. De Morgan sicher anwenden}
\textbf{a)} Bilde die Negation von $Y=A(B+\overline{C})$. Ziehe den Überstrich
Schritt für Schritt nach innen, bis kein Überstrich mehr über einer
Verknüpfung steht.
\antwort{26mm}
\textbf{b)} In einer Übung steht $\overline{A+BC}=\overline{A}+\overline{B}+\overline{C}$.
Begründe, warum das falsch ist, gib die richtige Form an und nenne eine
Belegung von $A$, $B$, $C$, für die sich beide Ausdrücke unterscheiden.
\antwort{30mm}
\textbf{c)} Ein NAND-Gatter soll als OR-Gatter mit zwei Eingangsblasen
gezeichnet werden. Formuliere in einem Satz, was beim Verschieben der
Ausgangsblase über den Gatterkörper mit dem Gattertyp und mit den Eingängen
geschieht.
\antwort{20mm}

\block{3. Die NAND-Aufbauten als Gleichung}
Es sei $\N(A,B)=\overline{AB}$. Gib für NOT, AND und OR jeweils die vollständige
Gleichungskette mit Zwischenwerten sowie die Zahl der NAND-Gatter an. Rechne
nach, dass sich die gewünschte Funktion ergibt, und benenne das Gesetz, das
den jeweiligen Schritt trägt.
\antwort{45mm}

\newpage
\titlepagepart{03 / NETZE RECHNEN · CA. 16 MINUTEN}{Von der Schaltung zur Funktion}

\block{4. Ein unbekanntes Netz bestimmen}
Gegeben ist das Netz $P=\N(A,A)$, \quad $Q=\N(A,B)$, \quad $Y=\N(P,Q)$.
\begin{center}
\begin{tikzpicture}[scale=0.85,transform shape]
\node[gate] (p) at (0,1) {NAND};\node[gate] (q) at (0,-1) {NAND};
\node[gate] (y) at (4,0) {NAND};
\draw[wire] (-2.8,1) node[left] {$A$}--(-1.7,1);\fill(-1.7,1)circle(1.5pt);
\draw[wire] (-1.7,1)|-($(p.west)+(0,.25)$);
\draw[wire] (-1.7,1)|-($(p.west)+(0,-.25)$);
\draw[wire] (-1.7,1)--(-1.7,-0.75)-|($(q.west)+(0,.25)$);
\draw[wire] (-2.8,-1.25) node[left] {$B$}--($(q.west)+(0,-.25)$);
\draw[wire] (p.east)--(2,1) node[above] {$P$}|-($(y.west)+(0,.25)$);
\draw[wire] (q.east)--(2,-1) node[below] {$Q$}|-($(y.west)+(0,-.25)$);
\draw[wire] (y.east)--++(1.2,0) node[right] {$Y$};
\end{tikzpicture}
\end{center}
\textbf{a)} Fülle die Tabelle vollständig aus.
\begin{center}
\begin{tabular}{cc|ccc}\toprule
$A$&$B$&$P$&$Q$&$Y$\\\midrule
0&0&\lin&\lin&\lin\\
0&1&\lin&\lin&\lin\\
1&0&\lin&\lin&\lin\\
1&1&\lin&\lin&\lin\\\bottomrule
\end{tabular}
\end{center}
\textbf{b)} Leite dasselbe Ergebnis algebraisch her, ohne die Tabelle zu
benutzen. Benenne die beiden Gesetze, die den Ausdruck zusammenfallen lassen.
\antwort{20mm}
\textbf{c)} Welche praktische Folgerung ziehst du daraus für eine Schaltung,
die du selbst gebaut hast?
\antwort{11mm}

\block{5. Halbaddierer: die Brücke zu Tag 4}
Für die Addition zweier einzelner Bits gilt
$S=A\oplus B$ und $C_{\text{out}}=AB$. Im XOR-Netz aus vier NAND-Gattern ist
$T=\N(A,B)$ der erste Zwischenwert.

\textbf{a)} Fülle die Tabelle aus und gib an, welche Binärzahl
$C_{\text{out}}S$ in der letzten Zeile darstellt.
\begin{center}
\begin{tabular}{cc|cc}\toprule
$A$&$B$&$S$&$C_{\text{out}}$\\\midrule
0&0&\lin&\lin\\
0&1&\lin&\lin\\
1&0&\lin&\lin\\
1&1&\lin&\lin\\\bottomrule
\end{tabular}
\end{center}
\textbf{b)} Zeige, dass sich $C_{\text{out}}$ allein aus dem vorhandenen $T$
gewinnen lässt. Gib die Gleichung an und nenne die Gesamtzahl der NAND-Gatter
einmal mit und einmal ohne Mitbenutzung von $T$.
\antwort{18mm}
\textbf{c)} Was ändert die Mitbenutzung von $T$ an der Wahrheitstabelle, und
was ändert sie an der realen Schaltung?
\antwort{13mm}

\newpage
\titlepagepart{04 / DIE ELEKTRISCHE SEITE · CA. 12 MINUTEN}{Pegel, Störabstand, Kennlinie}

\block{6. Störabstände berechnen}
Es gelten
\[
NM_L=V_{IL}-V_{OL},\qquad NM_H=V_{OH}-V_{IH}.
\]
\textbf{a)} Treiber: $V_{OL}=0{,}33$~V, $V_{OH}=2{,}97$~V. Empfänger:
$V_{IL}=0{,}99$~V, $V_{IH}=1{,}98$~V, $V_{DD}=3{,}3$~V. Berechne $NM_L$ und
$NM_H$ und beurteile, ob die Verbindung zulässig ist.
\antwort{22mm}
\textbf{b)} Zweite Paarung: Treiber $V_{OL}=0{,}50$~V, $V_{OH}=2{,}40$~V;
Empfänger $V_{IL}=0{,}40$~V, $V_{IH}=2{,}00$~V. Berechne beide Störabstände.
Einer wird negativ: Erkläre konkret, welche Fehlfunktion daraus folgt.
\antwort{24mm}
\textbf{c)} Ein Ausgang liefert HIGH mit $5{,}0$~V an einen Eingang einer
$3{,}3$-V-Familie. Der Wert liegt weit über $V_{IH}$. Warum ist die Verbindung
trotzdem unzulässig? Grenze logische von elektrischer Gültigkeit ab.
\antwort{22mm}

\block{7. Von der Kennlinie zur Regel}
\textbf{a)} Skizziere die DC-Übertragungskennlinie eines realen Inverters,
beschrifte beide Achsen und markiere die beiden Stellen, an denen
$\dfrac{dV_{out}}{dV_{in}}=-1$ gilt. Was wird an diesen Stellen festgelegt?
\antwort{40mm}
\textbf{b)} Formuliere die Static Discipline in einem Satz und nenne eine
Aussage, die sie ausdrücklich \emph{nicht} macht.
\antwort{20mm}
\textbf{c)} Welche der Teilaufgaben 6a--6c ließen sich in Logisim-Evolution
nachstellen, welche nicht? Begründe in zwei Sätzen.
\antwort{18mm}

\newpage
\titlepagepart{05 / PRAXIS UND ERFOLGSKONTROLLE · CA. 20 MINUTEN}{Die fünf Aufbauten abschließen}

\block{8. Die offenen NAND-Aufbauten}
Lege in Logisim-Evolution die Datei \code{tag05\_nand.circ} an und darin fünf
getrennte Teilschaltungen. Übernimm die bereits gebauten Netze, ergänze die
fehlenden. Alle Pins sind 1~Bit breit, alle NANDs haben zwei Eingänge.
\begin{center}
\begin{tabular}{cc|ccccc}\toprule
$A$&$B$&NOT ($Y=\overline{A}$)&AND&OR&XOR&NOR\\\midrule
0&0&\lin&\lin&\lin&\lin&\lin\\
0&1&\textemdash&\lin&\lin&\lin&\lin\\
1&0&\lin&\lin&\lin&\lin&\lin\\
1&1&\textemdash&\lin&\lin&\lin&\lin\\\bottomrule
\end{tabular}
\end{center}
NOT hat nur den Eingang $A$; die mit \textemdash{} markierten Felder entfallen.
Verwendete NAND-Anzahl: NOT \lin\ AND \lin\ OR \lin\ XOR \lin\ NOR \lin

\textbf{Kontrollfrage:} Weicht eine Zeile von der erwarteten Funktion ab,
welchen Zwischenwert prüfst du zuerst und warum?
\antwort{14mm}

\block{9. Die beiden De-Morgan-Nachweise protokollieren}
$Y_1$ ist die linke, $Y_2$ die rechte Seite der jeweiligen Gleichung.
\begin{center}
\begin{tabular}{cc|cc|cc}\toprule
&&\multicolumn{2}{c|}{$\overline{AB}$ gegen $\overline{A}+\overline{B}$}&\multicolumn{2}{c}{$\overline{A+B}$ gegen $\overline{A}\,\overline{B}$}\\
$A$&$B$&$Y_1$&$Y_2$&$Y_1$&$Y_2$\\\midrule
0&0&\lin&\lin&\lin&\lin\\
0&1&\lin&\lin&\lin&\lin\\
1&0&\lin&\lin&\lin&\lin\\
1&1&\lin&\lin&\lin&\lin\\\bottomrule
\end{tabular}
\end{center}
Warum ist erst die Übereinstimmung aller vier Zeilen ein Nachweis, eine
einzelne gleiche Zeile aber nicht?
\antwort{14mm}

\block{10. Erfolgskontrolle des Tages}
Der Plan verlangt fünf NAND-Schaltungen mit korrekter Wahrheitstabelle als
\code{.circ} im Repository. Hake ab und trage Fehlendes nach:
\begin{itemize}[leftmargin=8mm,itemsep=1pt,topsep=2pt]
\item[$\square$] \code{tag05\_nand.circ} enthält NOT, AND, OR, XOR und NOR, alle aus NAND.
\item[$\square$] Jede Teilschaltung wurde über alle Eingangskombinationen geprüft.
\item[$\square$] Beide De-Morgan-Vergleiche liegen als \code{.circ} vor.
\item[$\square$] \code{auswertung.md} hält die Tabellen dieser Seite und die Gatterzahlen fest.
\item[$\square$] \code{README.md} nennt die endgültigen Dateinamen.
\item[$\square$] Committet, etwa mit \code{git add 01\_digital} und
      \code{git commit -m "Tag 5: NAND-Schaltungen und De-Morgan-Nachweis"}.
\end{itemize}

\textbf{Abschlussnotiz für \code{LOG.md}:} Was sitzt nach diesem Tag, und was
hast du heute benutzt, ohne es zu verstehen? Zwei bis drei Sätze.
\antwort{16mm}

\end{document}
